\beginsong{Eu n-am să te las ca să pieri}
\beginverse
„Eu n-am să te las ca să pieri
Când treci prin necazul cel greu;
Atunci când Mă chemi în dureri,
Eu vin!”, a promis Dumnezeu!
\endverse

\beginchorus
/: „Cheamă-Mă în ziua cea grea,
Când nori de necaz vor veni,
Atunci Eu te voi ajuta,
Iar tu-n cânt Mă vei proslăvi! :/
\endchorus

\beginverse
Când ziua s-a dus, e târziu
Și singur în noapte rămâi,
În jurul tău totu-i pustiu
Și-o piatră de-ți pui căpătâi,
\endverse

\beginchorus
/: Cheamă-Mă și am să vin
Și-o scară la Cer voi ’nălța;
În noapte, sub cerul senin,
Cu tine de vorbă voi sta! :/
\endchorus

\beginverse
Te simți, de furtuni, clătinat
Și valul te-mpinge spre stânci,
În jur mai auzi, disperat,
Un vuiet de ape adânci...
\endverse

\beginchorus
/: Cheamă-Mă-n ziua cea grea
Când valul lovește cumplit;
În Mine, pe Stâncă, vei sta
Furtuna privind liniștit! :/
\endchorus

\beginverse
Eu sunt Cel ce sunt, Dumnezeu,
Și-aud când Mă strigi din necaz,
Căci Eu văd din cer chinul tău
Și lacrima pe-al tău obraz,
\endverse

\beginchorus
Crede și nu te-ndoi!
Așteaptă-n tăcere-ajutor,
Căci Eu lângă tine voi fi
Și chinul va fi mai ușor!
\endchorus

\beginchorus
Cheamă-Mă în ziua cea grea
Când nori de necaz vor veni;
Atunci Eu te voi ajuta,
Iar tu-n cânt Mă vei proslăvi!”
\endchorus
\endsong